\scanpage{32}
Un pion ne peut se placer en position de ``prise'', i.e.
sur une intersection libre telle qu'il n'ait plus
\inserted{aucune} \deleted{de} liberté (i.e. que le groupe qui le contient
n'ait plus de liberté) sauf dans les deux
cas suivants:

1°) Cas du ``Ko'' cf plus bas

2°) Cas de ``remplissage d'un trou'' : \deleted{disons que}
$A$ joue, sur une \deleted{\unsure{case}} \inserted{intersection $i$} qui est la seule liberté restante
d'un ensemble \inserted{$E$ réunion} de groupes de $B$ \inserted{(groupement de $B$)}, de façon que
toutes les libertés du groupe $G$ de $A$ contenant $i$
soient prises par des pions de $E$. ($E$ est entouré
par les pions de $A$ sauf en $i$, et $E$ entoure
le groupe \inserted{$G$} de $A$ contenant $i$)

\deleted{\unsure{Commençant à}}

\deleted{Chaque} \inserted{Les} joueurs prennent \deleted{\unsure{à tour de}} \inserted{à tour de} rôle, et ils ne
peuvent sauter leur tour (obligation de jouer).
L'obligation de jouer tombe dans le seul cas
suivant:

(fin de partie) Le joueur $A$ n'a plus de case
libre où il soit légitime de se mettre (\inserted{i.e.} où il
\unsure{n'en|n'y} soit pas en état de prise, sans \deleted{l'}\inserted{une des} exceptions
1° et 2° plus haut), sauf des \deleted{cases} \inserted{intersections $c$} du type
suivant: il y a un groupement $E$ de $A$
n'ayant plus que deux libertés \deleted{\unsure{intérieures}} \inserted{en $\unsure{i_1,i_2}$}
\deleted{\unsure{ces deux libertés}}
\mspage{31}
